\documentclass{ou}
\begin{document}
\thispagestyle{empty}
\begin{center}
\shortstack{\FGAVO}\vspace{4\baselineskip}
\textbf{РЕФЕРАТ}\\
\setstretch{2}
Информационные системы и~технологии\\
\setstretch{1}
Технологии создания и обработки аудио\vspace{2\baselineskip}\end{center}
\begin{flushleft}
\doublespacing
\begin{tabularx}{\textwidth}{@{} p{0.3\textwidth} l l @{}}
Руководитель & ст. преподаватель & И. Р. Нигматуллин\\
Студенты     &                   &                  \\
\multirow{-2.27}{\hsize}{\singlespacing Основы представления звука в цифровых системах}
             & ГФ25-01Б, 152537671 & А. Е. Корикова\\
             & ГФ25-01Б, 152536528 & Д. А. Васильева\\
             & ГФ25-01Б, 152540016 & И. Н. Морозова\\
\multirow{-2.27}{\hsize}{\singlespacing Методы кодирования, декодирования и~передачи аудиосигналов}
             & ГФ25-01Б, 152540046 & А. Д. Балцевич\\
             & ГФ25-01Б, 152537659 & В. Е. Зайцева\\
             & ГФ25-01Б, 152536592 & А. Ю. Бирюкова\\
\multirow{-2.27}{\hsize}{\singlespacing Отличие сжатия без~потерь и~с~потерями при~хранении аудио}
             & ГФ25-01Б, 152541294 & Э. Р. Махмудов\\
             & ГФ25-01Б, 152539829 & В. Д. Бобкова\\
\end{tabularx}
\end{flushleft}
\vspace*{\fill}
\begin{center}
\city
\end{center}
\newpage
\thispagestyle{empty}
\setstretch{1}
Продолжение титульного~листа реферата по~теме «Технологии создания и обработки аудио»
\begin{flushleft}
\doublespacing
\begin{tabularx}{\textwidth}{@{} p{0.3\textwidth} l l @{}}
Студенты     &                     &         \\
\multirow{-2.27}{\hsize}{\singlespacing Основные форматы звуковых файлов и~области их применения}
             & ГФ25-01Б, 152536115 & С. В. Идрисова\\
             & ГФ25-01Б, 152540131 & Е. Кичко\\
             & ГФ25-01Б, 152539526 & К. Н. Аркадьева\\
             & ГФ25-01Б, 152536056 & В. А. Карасова\\
\multirow{-2.27}{\hsize}{\singlespacing Основы редактирования и~монтажа аудио в~цифровой среде}
             & ГФ25-01Б, 152537684 & О. И. Бочанова\\
             & ГФ25-01Б, 152539986 & М. К. Гинтер\\
             & ГФ25-01Б, 152541471 & А. О. Безруких\\
             & ГФ25-01Б, 152536796 & А. В. Корюкова\\
\multirow{-2.27}{\hsize}{\singlespacing Перспективы развития технологий объемного и пространственного звука}
             & ГФ25-01Б, 152540765 & А. В. Гайворонская\\
             & ГФ25-01Б, 152540753 & М. И. Бобылев\\
             & ГФ25-01Б, 152539850 & Е. А. Аксаментова\\
            
\end{tabularx}
\end{flushleft}
\newpage
\tableofcontents
\newpage
\onehalfspacing
\section{Основы представления звука в цифровых системах}
Компьютерные устройства могут обрабатывать информацию в цифровой форме, которая является серией электрических импульсов, способных кодировать логические нули и единицы. Но звуковая информация, слышимая людьми, является непрерывной по своей природе. Это означает, что звуковые волны являются аналоговыми сигналами, имеющими переменную амплитуду и частоту. Для того чтобы записать такой звуковой сигнал в компьютерную память, нужно осуществить его преобразование в цифровую форму\cite{spravochnik}.

На сегодняшний день звуковые устройства стали неотъемлемой частью любого персонального компьютера. В процессе развития компьютерных технологий был сформирован стандарт звукового программно-аппаратного обеспечения. Звуковые устройства превратились из дорогих экзотических дополнений в привычный компонент компьютерных систем практически любых конфигураций.

Таким образом, фундаментальной задачей является преодоление разрыва между непрерывной природой звука и дискретной природой вычислительных систем.
\subsection{Аналоговая и цифровая формы сигнала}
В аналоговой звуковой аппаратуре аудиосигнал представляется в виде непрерывного электрического сигнала, параметры которого изменяются во~времени в соответствии с изменениями звукового давления. В отличие от этого, вычислительные системы оперируют дискретными данными. Следовательно, для обработки звука в компьютере требуется преобразование исходного аналогового сигнала в цифровую форму. Цифровой звук -- это дискретное представление аналогового электрического сигнала, определяемое последовательностью численных значений его амплитуды\cite{ixbt}. 

Цифровое представление электрических сигналов призвано внести в них избыточность, предохраняющую от воздействия паразитных помех. Для этого на несущий электрический сигнал накладываются серьезные ограничения -- его амплитуда может принимать только два предельных значения -- 0 и 1. Вся зона возможных амплитуд в этом случае делится на~три зоны: нижняя представляет нулевые значения, верхняя -- единичные, а промежуточная является запрещенной -- внутрь нее могут попадать только помехи. Таким образом, любая помеха, амплитуда которой меньше половины амплитуды несущего сигнала, не оказывает влияния на правильность передачи значений 0 и 1. Помехи с большей амплитудой также не оказывают влияния, если длительность импульса помехи ощутимо меньше длительности информационного импульса, а на входе приемника установлен фильтр импульсных помех.

Сформированный таким образом цифровой сигнал может переносить полезную информацию, которая закодирована в виде последовательности битов -- нулей и единиц; частным случаем такой информации являются электрические и звуковые сигналы. Здесь количество информации в~несущем цифровом сигнале значительно больше, нежели в~кодированном исходном, так что несущий сигнал имеет определенную избыточность относительно исходного, и любые искажения формы кривой несущего сигнала, при~которых еще сохраняется способность приемника правильно различать нули и~единицы, не влияют на достоверность передаваемой этим сигналом информации. Однако~в случае воздействия значительных помех форма сигнала может искажаться настолько, что точная передача переносимой информации становится невозможной -- в ней появляются ошибки, которые при простом способе кодирования приемник не сможет не только исправить, но и~обнаружить. Для еще большего повышения стойкости цифрового сигнала к помехам и~искажениям применяется цифровое избыточное кодирование двух типов: проверочные (EDC -- Error Detection Code, обнаруживающий ошибку код) и корректирующие (ECC -- Error Correction Code, исправляющий ошибку код) коды. Цифровое кодирование состоит в простом добавлении к исходной информации дополнительных битов и/или преобразовании исходной битовой цепочки в цепочку большей длины и другой структуры. EDC позволяет просто обнаружить факт ошибки -- искажение или выпадение полезной либо появление ложной цифры, однако переносимая информация в этом случае также искажается; ECC позволяет сразу же исправлять обнаруженные ошибки, сохраняя переносимую информацию неизменной. Для удобства и надежности передаваемую информацию разбивают на блоки (кадры), каждый из которых снабжается собственным набором этих кодов \cite{3dnews}.
\subsection{Теоретические предпосылки оцифровки звука}
Звуковая волна представляет собой сложную функцию, описывающую зависимость амплитуды от времени. Поскольку акустические колебания, как правило, не могут быть представлены аналитической функцией, их прямое математическое описание для последующего хранения в памяти ЭВМ невозможно.

Альтернативный метод заключается в дискретном описании функции путем фиксации ее значений в определенные моменты времени. Однако~данный подход имеет фундаментальное ограничение: значения амплитуды сигнала, являющиеся непрерывной величиной, не могут быть записаны с~абсолютной точностью и подлежат округлению. Таким образом, для~перехода в цифровую форму требуется аппроксимация исходного сигнала по двум осям: временной и амплитудной\cite{spravochnik}.
\subsection{Аналогово-цифровой и цифро-аналоговый преобразователи}
Первый преобразует аналоговый сигнал в цифровое значение амплитуды, второй выполняет обратное преобразование. 

Принцип работы АЦП~состоит~в измерении уровня~входного~сигнала и~выдаче результата~в~цифровой~форме. Непрерывный~аналоговый~сигнал превращается в~импульсный, с~одновременным измерением амплитуды каждого импульса. ЦАП получает на входе цифровое значение амплитуды и выдает на выходе импульсы напряжения или тока нужной величины, которые расположенный за ним интегратор (аналоговый фильтр) превращает в непрерывный аналоговый сигнал.

Для правильной работы АЦП входной сигнал не должен изменяться в~течение времени преобразования, для чего на его входе обычно помещается схема выборки-хранения, фиксирующая мгновенный уровень сигнала и~сохраняющая его в течение всего времени преобразования. На выходе ЦАП также может устанавливаться подобная схема, подавляющая влияние переходных процессов внутри ЦАП на параметры выходного сигнала\cite{ixbt}.
\subsection{Преобразование аналогового сигнала в цифровую форму}
Этапы аналого-цифрового преобразования:
\begin{itemize}[label=-, leftmargin=0em, itemindent=3.35em, nosep]
\item Дискретизация;
\item Квантование.
\end{itemize}

Процесс дискретизации заключается в получении значений амплитуды преобразуемого сигнала через равные промежутки времени. Количество таких измерений, производимых за одну секунду, называется частотой дискретизации. Данный процесс позволяет перейти от непрерывной функции к последовательности ее дискретных значений.

Процесс квантования заключается в замене реальных, непрерывных значений амплитуды на приближенные, принадлежащие к конечному набору дискретных уровней. Точность этого процесса определяется разрядностью квантования, то есть количеством бит, выделяемых для кодирования одного отсчета. Чем выше разрядность, тем большее количество уровней квантования доступно и тем меньше ошибка округления\cite{spravochnik}.
\subsection{Теорема Котельникова и ее практическое применение}
Согласно теореме Котельникова, любой непрерывный процесс с~ограниченным спектром может быть полностью описан дискретной последовательностью его мгновенных значений, следующих с частотой, как~минимум вдвое превышающей частоту наивысшей гармоники процесса; частота Fd выборки мгновенных значений (отсчетов) называется частотой дискретизации.

Из теоремы следует, что сигнал с частотой Fa может быть успешно дискретизирован по времени на частоте 2Fa только в том случае, если он является чистой синусоидой, ибо любое отклонение от синусоидальной формы приводит к выходу спектра за пределы частоты Fa. Таким образом, для временной дискретизации произвольного звукового сигнала (обычно имеющего, как известно, плавно спадающий спектр), необходим либо выбор частоты дискретизации с запасом, либо принудительное ограничение спектра входного сигнала ниже половины частоты дискретизации.

Одновременно с временной дискретизацией выполняется амплитудная -- измерение мгновенных значений амплитуды и их представление в виде числовых величин с определенной точностью. Точность измерения определяет соотношение сигнал/шум и динамический диапазон сигнала\cite{ixbt}.
\subsection{Импульсно-кодовая модуляция}
Полученный поток чисел, описывающий звуковой сигнал, называют импульсно-кодовой модуляцией или ИКМ (Pulse Code Modulation, PCM), так~как каждый импульс дискретизованного по~времени сигнала представляется собственным цифровым кодом. PCM-поток может быть как~параллельным, когда все биты каждого отсчета передаются одновременно по нескольким линиям с частотой дискретизации, так~и~последовательным, когда биты передаются друг за другом с более высокой частотой по одной линии.

Чаще всего применяют линейное квантование, когда числовое значение отсчета пропорционально амплитуде сигнала. Из-за логарифмической природы слуха более целесообразным было бы логарифмическое квантование, когда числовое значение пропорционально величине сигнала в децибелах, однако~это сопряжено с трудностями чисто технического характера\cite{3dnews}.
\subsection{Стандарты качества}
В большинстве современных цифровых звуковых систем используются стандартные частоты дискретизации 44,1 и 48 кГц, однако частотный диапазон сигнала обычно ограничивается возле 20 кГц для оставления запаса по~отношению к теоретическому пределу. Также наиболее распространено 16-разрядное квантование по уровню, что дает предельное соотношение сигнал/шум около 98 дБ. В студийной аппаратуре используются более высокие разрешения -- 18-, 20- и 24-разрядное квантование при частотах дискретизации 56, 96 и 192 кГц. Это делается для того, чтобы сохранить высшие гармоники звукового сигнала, которые непосредственно не воспринимаются слухом, но~влияют на формирование общей звуковой картины\cite{ixbt}.
\subsection{Преимущества и недостатки цифрового звука}
Цифровое представление звука ценно возможностью бесконечного хранения и тиражирования без потери качества, но~преобразование из~аналоговой формы в цифровую и обратно все же неизбежно приводит к~частичной его потере. Наиболее неприятные на слух искажения, вносимые на этапе оцифровки -- гранулярный шум, возникающий при квантовании сигнала по уровню из-за округления амплитуды до ближайшего дискретного значения. В отличие от простого широкополосного шума, вносимого ошибками квантования, гранулярный шум представляет собой гармонические искажения сигнала, наиболее заметные в верхней части спектра.

Искажения, вносимые гранулярным шумом, можно уменьшить путем добавления к сигналу обычного белого. Это приводит к незначительному увеличению уровня шума, зато ослабляет корреляцию ошибок квантования с высокочастотными компонентами сигнала и улучшает субъективное восприятие. Сглаживание применяется также перед округлением отсчетов при~уменьшении их разрядности. 

При восстановлении звука из цифровой формы в аналоговую появляется проблема сглаживания ступенчатой формы сигнала, подавления гармоник, вносимых частотой дискретизации. Из-за неидеальности амплитудной характеристики фильтров может происходить либо недостаточное подавление этих помех, либо избыточное ослабление полезных высокочастотных составляющих. Плохо подавленные гармоники частоты дискретизации искажают форму аналогового сигнала (особенно в области высоких частот), что создает впечатление «шероховатого», «грязного» звука\cite{3dnews}.
\newpage
\section{Методы кодирования, декодирования и~передачи аудиосигналов}
Методы кодирования, декодирования и передачи аудиосигналов лежат в основе современных технологий хранения, обработки и передачи звука. Аудиосигнал необходимо преобразовать в цифровую форму, чтобы его можно было анализировать, передавать и сжимать. 

Кодирование звуковой информации включает оцифровку аналогового сигнала и различные методы сжатия, которые можно разделить на сжатие с~потерями и сжатие без потерь.
\subsection{Оцифровка аудиосигнала}
Цифровой звук -- аналоговый звуковой сигнал, представленный посредством дискретных численных значений его амплитуды.

Оцифровка звука -- технология осуществления замеров амплитуды звукового сигнала с определенным временным шагом и последующей записи полученных значений в численном виде. Другое название оцифровки звука -- аналогово-цифровое преобразование звука.

Оцифровка преобразует аналоговый звуковой сигнал в цифровой. Она включает два этапа:
\begin{enumerate}[leftmargin=0em, itemindent=3.75em, nosep, label=\arabic*)]
\item Дискретизация;
\item Квантование.
\end{enumerate}

Дискретизация -- измерение амплитуды звука через равные промежутки времени, процесс получения значений сигнала, который преобразуется с~определенным временным шагом -- шагом дискретизации. Количество замеров величины сигнала, осуществляемых в единицу времени, называют частотой дискретизации. Чем меньше шаг дискретизации, тем выше частота дискретизации и тем более точное представление о сигнале нами будет получено.

Квантование -- округление измеренного значения до ближайшего уровня. Результатом является цифровое представление звука, например PCM. Точность округления зависит от выбранного количества (2N) уровней квантования, которое, в свою очередь, зависит от количества бит (N), отведенных для~записи значения амплитуды. Принимается, что погрешности квантования, являющиеся результатом квантования с разрядностью 16 бит, остаются для слушателя почти незаметными.
\subsection{Кодирование аудиосигнала}
Цель кодирования -- уменьшить размер данных, облегчить передачу и~сохранить нужное качество.

Методы кодирования:
\begin{itemize}[label=-, leftmargin=0em, itemindent=3.35em, nosep]
\item Кодирование без потерь;
\item Кодирование с потерями.
\end{itemize}

Кодирование без потерь позволяет полностью восстановить исходный сигнал. Использует предсказание сигнала и энтропийное кодирование. Применяется в архивировании и профессиональной звукозаписи. Примеры аудиоформатов: FLAC, ALAC

Кодирование с потерями удаляет компоненты, которые человеческое ухо не воспринимает. Значительно снижает объём данных. Используется в~стриминге, мобильных устройствах, онлайн-видео. Примеры аудиоформатов: MP3, AAC.

Декодирование аудиосигнала -- это процесс восстановления звука из~закодированного цифрового представления.

Если кодирование уменьшает размер данных или подготавливает их для передачи, то декодирование возвращает структуру сигнала к состоянию, пригодному для воспроизведения.

Результатом является цифровой или аналоговый звук, пригодный для~воспроизведения.

Основные методы декодирования:
\begin{enumerate}[leftmargin=0em, itemindent=3.75em, nosep, label=\arabic*)]
\item Психоакустическое декодирование. Используется в форматах с~потерями (MP3, AAC), основано на восстановлении сигнала с учётом особенностей слуха человека;
\item Декодирование на основе частотных преобразований. Используется в форматах, где кодирование работало в спектральной области, примеры форматов: MP3, AAC, OGG Vorbis, Opus (частично);
\item Декодирование параметрического аудио. Метод используется в~кодеках низкого битрейта. Передаются только параметры звука (тон, форманты, шум, энергетика, частота), а декодер синтезирует звук;
\item Декодирование на основе временных моделей. Временной сигнал восстанавливается напрямую, без спектральных операций. Используется в~некоторых простых кодеках или в режиме SILK (Opus). При этом декодер восстанавливает форму сигнала по предсказанию и разнице (residual error).
\end{enumerate}
\subsection{Передача аудиосигналов}
Передача аудиосигналов -- это процесс доставки звуковой информации от~источника к приемнику через проводные или беспроводные каналы связи.

Передача может быть локальной или сетевой и включает:
\begin{itemize}[label=-, leftmargin=0em, itemindent=3.35em, nosep]
\item Компрессию -- уменьшение размера данных перед отправкой;
\item Мультиплексирование -- объединение нескольких потоков;
\item Коррекцию ошибок -- защита от искажений;
\item Сетевые протоколы -- управление передачей данных.
\end{itemize}

Примеры используемых кодеков:
\begin{enumerate}[leftmargin=0em, itemindent=3.75em, nosep, label=\arabic*)]
\item Для музыки -- AAC, MP3, Opus;
\item Для речи -- AMR, Speex;
\item Для профессионального звука -- PCM, FLAC.
\end{enumerate}

Методы передачи аудиосигналов:
\begin{itemize}[label=-, leftmargin=0em, itemindent=3.35em, nosep]
\item Аналоговые методы передачи. Используются в традиционном радио и~старых системах связи. Передается непрерывный сигнал;
\item Цифровые методы передачи. Сигнал предварительно преобразуется в~набор битов. Это основной современный подход;
\item Цифровая модуляция. Нужна, чтобы передать цифровые данные по~физическому каналу (радио, кабель);
\item Передача аудио через интернет.
\end{itemize}

Методы кодирования и передачи аудиосигналов используются в~телефонии, интернет-связи, стриминговых сервисах, телевидении и радио, профессиональных студиях и мультимедийных системах.
\newpage
\section{Отличие сжатия без~потерь и~с~потерями при~хранении аудио}
Сжатие аудиоданных -- это процесс уменьшения скорости цифрового потока за счет устранения избыточности, содержащейся в исходном сигнале. Несмотря на существование несжатых форматов (таких как WAV, AIFF), для эффективного хранения и передачи звука повсеместно используются алгоритмы сжатия, реализуемые кодеками. Все методы сжатия аудио можно разделить на два принципиально различных класса: сжатие без потерь (lossless) и сжатие с потерями (lossy).
\subsection{Сжатие без потерь} 
Данный подход, представленный такими форматами, как FLAC и ALAC, ставит своей целью полное восстановление исходных данных из сжатого представления. В его основе лежит устранение статистической избыточности цифрового звукового сигнала.

Основной принцип заключается в учете корреляционной связи между соседними отсчетами сигнала. Устранение этой избыточности позволяет сократить объем данных на 15–25\verb|%|. Для дальнейшего уменьшения скорости цифрового потока применяются методы энтропийного кодирования, учитывающие статистику сигнала. Ярким примером является код Хаффмана, в котором наиболее вероятным значениям сигнала присваиваются короткие кодовые слова, а маловероятным -- более длинные. Таким образом, сжатие без потерь по своим методам близко к алгоритмам архивации данных и~гарантирует битовую идентичность восстановленного файла оригиналу, хотя~и обеспечивает относительно скромную степень сжатия.
\subsection{Сжатие с потерями} 
Этот подход, используемый в форматах MP3, AAC и других, обеспечивает значительно более высокую степень сжатия (в 10–12 раз) за счет необратимого отбрасывания части информации. Восстановление данных происходит с~искажениями, однако эти искажения проектируются таким образом, чтобы~быть минимально заметными для человеческого слуха. В основе лежит устранение психоакустической избыточности.
\subsection{Психоакустическая модель и эффект маскирования}
Алгоритмы сжатия с потерями опираются на несовершенство человеческого слухового восприятия, в частности, на эффекты слухового маскирования.

Частотное (одновременное) маскирование: в присутствии громкого звука~-- «маскера» на определенной частоте, более тихие звуки на близлежащих частотах становятся неслышимыми. Кривая порога маскирования показывает, как присутствие маскера повышает абсолютный порог слышимости в его частотном окружении.

Временное (неодновременное) маскирование: громкий звук может маскировать более тихие звуки, которые возникают непосредственно до~или~после него, используя инерционность слуховой системы.

\subsection{Техническая реализация}
На практике эти принципы реализуются с помощью полосного кодирования:
\begin{itemize}[label=-, leftmargin=0em, itemindent=3.35em, nosep]
\item Входной сигнал разбивается на кадры и проходит через банк фильтров, разделяющий его на частотные поддиапазоны;
\item Для каждого поддиапазона психоакустическая модель анализирует спектральный состав и вычисляет пороги маскирования, определяя, какая часть сигнала может быть отброшена или переквантована с меньшим числом бит без~заметной потери качества;
\item Производится масштабирование и перераспределение битов между полосами, после чего данные кодируются.
\end{itemize}

Поскольку в процессе сжатия удаляется психоакустическая избыточность, точное восстановление исходного сигнала невозможно.

\subsection{Сравнительный анализ и область применения}

Ключевое различие между методами заключается в результате декодирования: lossless-кодеки обеспечивают полное восстановление данных, в то время как lossy-кодеки -- лишь приближенное.


Сжатие без потерь применяется в ситуациях, где критически важна сохранность битовой точности аудиозаписи: в профессиональной звукозаписи, архивации фонограмм, а также аудиофилами. Его недостаток -- больший объем получаемых файлов по сравнению с lossy-форматами.

Сжатие с потерями доминирует в областях, где приоритетом является эффективное использование пропускной способности каналов и дискового пространства: в потоковой передаче данных (стриминговые сервисы), персональных аудиоколлекциях и цифровой дистрибуции. Эффективность сжатия позволяет значительно экономить трафик и место для хранения.
\newpage
\section{Основные форматы звуковых файлов и~области их применения}
Для каждого, кто хоть немного связан с миром музыки, не секрет, что цифровые аудиофайлы могут храниться в различных форматах. Однако~не~все понимают, чем эти форматы отличаются друг от друга. Знание данной особенности важно для эффективного использования сервисов, связанных с~созданием, обработкой и хранением музыки на разных носителях.
\subsection{Понимание процесса оцифровки звука}
Прежде чем углубиться в мир аудиоформатов, стоит разобраться с~тем, как происходит процесс оцифровки звука. Он реализуется с помощью аналого-цифрового преобразователя (АЦП), который преобразует аналоговый звуковой сигнал в цифровой формат, применяя метод импульсного кодирования. АЦП фиксирует изменения амплитуды сигнала, создавая цифровую копию аналогового звука. Процедуру оцифровки часто называют дискретизацией -- от латинского слова discretus, что означает «прерывистый». Важный вопрос, который при этом возникает: какая частота дискретизации позволяет максимально точно передать аналоговый звук в цифровом формате?

Частота дискретизации -- это количество измерений входного сигнала в секунду, измеряемое в герцах (Гц). Одно измерение соответствует частоте в 1 Гц, тысяча измерений в секунду -- 1 килогерц (кГц). Порог слышимости для человеческого уха составляет 20 кГц, и чтобы оцифрованный звук не~отличался от аналогового, дискретизация должна происходить минимум 40~000 раз в~секунду.
\subsection{Основные параметры цифрового звука: частота дискретизации и битрейт}
В музыкальной индустрии особое значение имеют две частоты дискретизации:
\begin{enumerate}[leftmargin=0em, itemindent=3.75em, nosep, label=\arabic*)]
\item 44.1 кГц -- стандарт для CD и большинства цифровых записей. Данный показатель был выбран исторически, исходя из особенностей использования магнитной видеопленки в начале 80-х годов. Пленка записывала 60 кадров в~секунду, каждый из которых содержал 24 строки выборки. Перемножив эти значения, получаем 44100 выборок -- отсюда и частота дискретизации;
\item 48 кГц -- этот стандарт пришел из киноиндустрии, где пленки записывались с частотой 24 кадра в секунду, а синхронизация со звуком давала частоту 48 кГц. Частоты 88.2 кГц, 96 кГц и 192 кГц используются в~профессиональной звукозаписи для получения максимально высокого качества аудио.
\end{enumerate}
\subsection{Три основных категории аудиоформатов}
В мире цифровой музыки аудиоформаты делятся на три категории: форматы без сжатия, форматы со сжатием без потерь и форматы со сжатием с~потерями. Рассмотрим их подробнее.

Форматы без сжатия:
\begin{enumerate}[leftmargin=0em, itemindent=3.75em, nosep, label=\arabic*)]
\item WAV/WAVE (Waveform Audio File Format);
\item AIFF.
\end{enumerate}

WAV/WAVE -- это формат, появившийся в 1991 году для записи на~компакт-диски с частотой дискретизации 44.1 кГц и 16-битной разрядностью. Разработчиком формата выступили компании Microsoft и IBM в рамках разработки спецификации RIFF (Resource Interchange File Format) для~хранения мультимедийных данных на ПК. WAV сохраняет звук без каких-либо изменений. В связи с этим он часто используется для записи исходных музыкальных файлов. Одна минута аудио в формате WAV занимает 4–5 Мб в зависимости от частоты дискретизации и разрядности. На самом деле, стандартная минута стереозаписи с параметрами CD (44.1 кГц/16 бит) занимает около 10 МБ. Этот формат является основным для профессиональной звукозаписи, сведения и мастеринга в студиях, так как гарантирует работу с~чистыми, неизмененными данными.

AIFF создан компанией Apple в 1988 году. Является функциональным аналогом WAV для экосистемы Apple. Как и WAV, сохраняет звук без потерь, что делает его идеальным для аудиомонтажа и хранения мастер-записей в~профессиональных музыкальных и видеоредакторах.

Форматы со сжатием без потерь:
\begin{enumerate}[leftmargin=0em, itemindent=3.75em, nosep, label=\arabic*)]
\item FLAC (Free Lossless Audio Codec);
\item ALAC (Apple Lossless Audio Codec);
\item APE (Monkey's Audio).
\end{enumerate}

FLAC -- один из самых популярных форматов, который появился в~начале 2000-х. Разработчиком является некоммерческая организация Xiph.Org Foundation. Он позволяет сжать музыкальные файлы в 1.5–3.5 раза, сохраняя при этом все качество исходного звука. Частота дискретизации может достигать 192 кГц при разрядности 24 бит. Будучи открытым и бесплатным форматом, FLAC широко используется аудиофилами для архивирования музыкальных коллекций и распространения музыки в высоком разрешении, так как он обеспечивает качество, идентичное оригиналу, но при меньшем размере файла.

ALAC -- собственный формат Apple, первоначально разработанный компанией Apple в 2004 году, а в 2011 году ставший открытым и бесплатным. ALAC поддерживает частоты до 384 кГц и разрядность до 32 бит. ALAC идеально воспроизводится на устройствах Apple благодаря встроенным музыкальным процессорам DSP. Он является полноценной альтернативой FLAC для пользователей iPhone, iPad и Mac, позволяя им экономить место на~устройстве без потери качества звука.

APE -- формат для Windows, разработанный Мэтью Т. Эшлендом. Monkey's Audio поддерживает 8-, 16- и 24-битные файлы. Его широко используют в профессиональной аудиозаписи, хотя поддержка на других платформах ограничена. Несмотря на высокую степень сжатия, он требует значительных вычислительных ресурсов для кодирования и декодирования, поэтому его популярность ниже, чем у FLAC.

Форматы со сжатием с потерями:
\begin{enumerate}[leftmargin=0em, itemindent=3.75em, nosep, label=\arabic*)]
\item MP3 (MPEG-1 Audio Layer 3);
\item AAC (Advanced Audio Coding);
\item WMA;
\item OGG VORBIS.
\end{enumerate}

MP3 -- самый распространенный формат сжатия, который позволяет уменьшить размер файла до 11 раз за счет удаления неслышимых человеческим ухом частот. Разработка формата велась в конце 1980-х годов в немецком Институте Фраунгофера под руководством Карлхайнца Бранденбурга. MP3 файлы могут иметь битрейт от 32 до 320 кбит/с и частоту дискретизации до 48~кГц. Алгоритм MP3 использует психоакустическую модель, отсекая те звуки, которые маскируются более громкими и яркими частотами. Это делает его идеальным для повседневного прослушивания и хранения больших библиотек на портативных устройствах.

AAC создан консорциумом компаний включая Fraunhofer IIS, Dolby, Sony и Nokia. Обеспечивает лучшее качество, чем MP3, при одинаковом битрейте. Является стандартом для YouTube, iTunes и многих стриминговых сервисов.

WMA -- формат от Microsoft, который обеспечивает хорошее сжатие и~качество звука, хотя его поддержка ограничена в основном устройствами на~базе Windows. Исторически он создавался как конкурент MP3, но широкого распространения за пределами экосистемы Microsoft не получил.

OGG VORBIS -- некоммерческий формат с открытым исходным кодом, который обеспечивает лучшее качество звука по сравнению с MP3 при аналогичном битрейте. Данная разновидность используется в играх и различных приложениях, хотя его популярность ниже. Благодаря своей открытости и отсутствию лицензионных отчислений, Ogg Vorbis является популярным выбором для разработчиков игр (в таких движках, как Unity и~Unreal) и программного обеспечения с открытым исходным кодом.

Помимо частоты дискретизации, важную роль играет и битрейт. Данным термином называют количество данных, уходящих на кодирование одной секунды звука. Битрейт рассчитывается как произведение частоты дискретизации на разрядность и количество каналов. Разрядность показывает, сколько бит используется для кодирования каждого изменения звука, и~напрямую влияет на точность записи. Для стереозаписи на CD (2 канала): 44,100 Гц × 16 бит × 2 = 1,411,200 бит/с = 1411 кбит/с. В форматах со сжатием битрейт становится главным показателем качества: чем он выше, тем меньше данных было отброшено при сжатии.
\subsection{Значение понимания различия в аудиоформатах}
Для обычного слушателя выбор аудиоформата может не казаться важным, но меломаны и профессионалы должны уметь различать их, чтобы получать максимальное удовольствие от прослушивания. Лучшие аудиоформаты с~высоким битрейтом лучше раскрываются на специализированной технике, представленной CD-проигрывателями, системой домашних кинотеатров и~аудиопроигрывателями высокого класса.
\newpage
\section{Основы редактирования и~монтажа аудио в~цифровой среде}
\newpage
\section{Перспективы развития технологий объемного и~пространственного звука}
\subsection{Понятие объёмного и пространственного звука}
Пространственный или объёмный звук -- это технология, позволяющая воспроизводить аудиосреду так, чтобы слушатель ощущал направление, глубину и локализацию источников звука. В отличие от стерео, который создаёт иллюзию ширины сцены, объёмный звук формирует трёхмерное аудиопространство.

Основные характеристики:
\begin{itemize}[label=-, leftmargin=0em, itemindent=3.35em, nosep]
\item Локализация звука -- способность определить направление источника;
\item Глубина сцены -- восприятие расстояния;
\item Иммерсивность -- эффект присутствия;
\item Объёмность -- распределение звуков вокруг слушателя во всех направлениях.
\end{itemize}

Пример применения: в кинотеатрах формата Dolby Atmos установка потолочных и боковых динамиков позволяет зрителю услышать перемещение объектов (например, полёт самолёта или прошедший мимо поезд) так, будто они находятся в реальном пространстве.
\subsection{Современные технологии трёхмерного звука}
Dolby Atmos -- одна из самых популярных технологий пространственного звука. Использует объектно-ориентированную модель звука, где звук рассматривается как объект, помещённый в 3D-пространство.

Преимущества:
\begin{itemize}[label=-, leftmargin=0em, itemindent=3.35em, nosep]
\item Точная локализация аудиообъектов;
\item Поддержка кинотеатров, домашнего аудио и мобильных устройств;
\item Масштабируемость под разные конфигурации динамиков.
\end{itemize}

Пример: современные смартфоны поддерживают Dolby Atmos в фильмах и играх, создавая эффект объёмного звучания даже без дополнительных устройств.

DTS:X -- альтернативная технология, также использующая объектный звук. DTS:X гибко адаптируется под доступные динамики и позволяет регулировать громкость отдельных звуков (например, голосов актёров). Пример: домашние кинотеатры премиум-класса предлагают систему DTS:X, что делает звук более реалистичным и настраиваемым под конкретное помещение.

Sony 360 Reality Audio -- технология ориентирована на музыкальные сервисы и наушники. Каждый инструмент в записи размещается в виртуальной сфере вокруг слушателя.

Возможности:
\begin{itemize}[label=-, leftmargin=0em, itemindent=3.35em, nosep]
\item Иммерсивные музыкальные концерты;
\item Создание студийных треков с 360-градусным звучанием.
\end{itemize}

Spatial Audio (Apple). Apple внедрила пространственный звук в AirPods и~своих устройствах, используя отслеживание движения головы (head tracking).

Преимущества:
\begin{itemize}[label=-, leftmargin=0em, itemindent=3.35em, nosep]
\item Адаптация звуковой сцены под повороты головы;
\item Интеграция с VR/AR;
\item Реалистичные голосовые звонки и FaceTime-конференции.
\end{itemize}

Binaural Audio. Бинауральный звук создаётся через моделирование работы человеческих ушей. Он используется в VR и ASMR-контенте. Пример: В VR-играх бинауральный звук помогает игроку слышать приближающиеся шаги врагов с высокой точностью.
\subsection{Тенденции развития технологий пространственного звука}
Традиционные каналы больше не ограничивают разработчиков. Объектный звук позволяет:
\begin{itemize}[label=-, leftmargin=0em, itemindent=3.35em, nosep]
\item Не привязываться к числу динамиков;
\item Более точно воспроизводить перемещение объектов;
\item Передавать сложные сцены (метель, толпа, дождь) реалистично.
\end{itemize}

Также персонализация учитывает анатомические особенности ушей пользователя (HRTF -- Head-Related Transfer Function). Перспективы:
\begin{enumerate}[leftmargin=0em, itemindent=3.75em, nosep, label=\arabic*)]
\item Индивидуальные настройки для каждого пользователя;
\item Кастомизация звучания в наушниках;
\item Улучшенная локализация.
\end{enumerate}

Пример: некоторые сервисы анализируют фотографию уха пользователя, чтобы настроить пространственное звучание.

Интеграция искусственного интеллекта помогает:
\begin{itemize}[label=-, leftmargin=0em, itemindent=3.35em, nosep]
\item Реконструировать звуковую сцену;
\item Удалять шумы;
\item Создавать пространственный звук из обычного стерео.
\end{itemize}

Пример: платформы потоковой музыки используют ИИ для «апмикса» (повышения) стереотреков до формата 3D-звука без участия звукорежиссёра.

VR/AR становится главным драйвером развития пространственного звучания, так как зрительная иммерсия невозможна без звуковой.

Тенденции:
\begin{enumerate}[leftmargin=0em, itemindent=3.75em, nosep, label=\arabic*)]
\item Реалистичное звуковое позиционирование;
\item Синхронизация движений головы и аудиосцены;
\item Моделирование акустики помещений в реальном времени.
\end{enumerate}

Пример: VR-симуляторы для пилотов широко используют технологии пространственного звука для реалистичного восприятия приборов и внешней обстановки.

Активное развитие получают:
\begin{itemize}[label=-, leftmargin=0em, itemindent=3.35em, nosep]
\item Наушники с пространственным звуком;
\item Компактные саундбары, имитирующие многоканальную акустику;
\item Портативные колонки с 360-градусным воспроизведением.
\end{itemize}
\subsection{Перспективные направления развития}
Будущие системы смогут:
\begin{enumerate}[leftmargin=0em, itemindent=3.75em, nosep, label=\arabic*)]
\item Моделировать акустику любых помещений (от концертного зала до~небольшого офиса);
\item Автоматически подстраивать звук под помещение;
\item Создавать полностью виртуальные звуковые среды.
\end{enumerate}

Пример: музыкант сможет записывать трек в студии, затем воспроизвести его «как будто» он звучал в соборе или на стадионе.

Метавселенные требуют реалистичного трёхмерного звука. Перспективы:
\begin{itemize}[label=-, leftmargin=0em, itemindent=3.35em, nosep]
\item Многопользовательские концерты в VR;
\item Виртуальные офисы с индивидуальными зонами слышимости;
\item Звуковая навигация в 3D-пространстве.
\end{itemize}

Голографический звук (Holophonic audio) -- технология будущего, основанная на волновом поле. Преимущества:
\begin{enumerate}[leftmargin=0em, itemindent=3.75em, nosep, label=\arabic*)]
\item Абсолютная точность локализации;
\item Возможность «передвигать» источники звука в комнате;
\item Создание аудио-голограмм.
\end{enumerate}

Перспективы применения биометрии и нейротехнологий:
\begin{itemize}[label=-, leftmargin=0em, itemindent=3.35em, nosep]
\item Адаптация звука по реакции мозга;
\item Создание «идеального» звука на основе эмоций;
\item Управление аудиосценой с помощью интерфейсов.
\end{itemize}

Слияние пространственного звука и искусственного зрения. Будущие устройства смогут:
\begin{enumerate}[leftmargin=0em, itemindent=3.75em, nosep, label=\arabic*)]
\item Отслеживать положение тела и головы;
\item Анализировать окружающее пространство;
\item Создавать звуковую сцену, соответствующую реальной обстановке.
\end{enumerate}

Пример: система в автомобиле может предупреждать водителя о~приближающемся объекте, при этом создавая ощущение присутствия в~соответствующем направлении.
\subsection{Практические возможности применения новых технологий}
Киноиндустрия. Реалистичные спецэффекты, объёмная музыка, улучшенный опыт IMAX и VR-кино.

Видеоигры. Игра -- одна из главных сфер, где пространственный звук значительно повышает точность и вовлечённость. Пример: в шутерах игроки слышат шаги врагов с точностью до метра.

Медицина. Использование 3D-звука для:
\begin{enumerate}[leftmargin=0em, itemindent=3.75em, nosep, label=\arabic*)]
\item Хирургических симуляторов;
\item Реабилитации пациентов с нарушениями слуха;
\item Диагностики.
\end{enumerate}

Смарт-устройства и бытовая техника. Умные колонки, телевизоры, саундбары всё чаще используют алгоритмы пространственного звука.

Музыка будущего. Будут создаваться специальные треки с трёхмерным звучанием, а концерты смогут проходить полностью в виртуальном пространстве.
\newpage
{
\sloppy
\nocite{*}
\cleardoublepage
\addcontentsline{toc}{section}{Список использованных источников}
\bibliographystyle{plainurl}
\bibliography{report8}
}
\end{document}
